\chapter{Introduction}

\section{Motivation}
\begin{description}
\item[Efficiency]  I did not want conversions between textual and native
  representations for arithmetic types.  For example, if I need to stream
  an unsigned 32-bit integer of value 0xFFFFFFFF, I should only need to
  transfer 4 bytes, not 10 or more.
\item[Portability]  Since network transport was a primary goal, I was
  motivated to provide schemes that work across fairly common hardware
  platforms and UNIX operating systems.
\item[Flexibility]  I/O facilities should be flexible enough to accomodate
  instances of common class templates from the standard library, particularly
  containers like \texttt{vector}, \texttt{map}, et. al.  They should
  also be flexible enough to accomodate user-defined types without a
  tremendous amount of work.
\end{description}

\section{History}
libDwm was started around 1998.  Needless to say, C++ was a very
different language in 1998 versus 2025.  And the standard library has
grown tremendously since 1998.

The I/O facilities in libDwm started appearing in the early 2000's
(around the time we got C++03).  Since that time, it has evolved a bit
with each new version of the C++ standard, and is about to see a sea
change with C++26 reflection (P2996 and P3394 in particular).

\section{Primary I/O classes}

\section{Directly supported types}

\section{Supported \texttt{std} class templates}

\section{Concepts}
